\section{从 CLIP 引导到扩散模型引导}

\subsection{CLIP 引导的 3D 生成}

在 Score Distillation 出现之前，研究者首先探索了使用 CLIP（Contrastive Language-Image Pre-training）模型的先验知识来辅助 3D 重建。

\begin{knowledgebox}{CLIP 模型简介}
CLIP 是 OpenAI 提出的视觉-语言对齐模型，它将图像和文本映射到同一个嵌入空间中。通过计算嵌入向量之间的距离，可以衡量一张图像与一段文本描述之间的语义对齐程度。CLIP 在大规模图文对数据上训练，具有强大的零样本迁移能力。
\end{knowledgebox}

\subsubsection{少样本 3D 重建}

利用 CLIP 的思路如下：
\begin{itemize}
    \item 对于有参考图像的视角：使用传统的 L2 重建损失
    \item 对于没有参考图像的视角：使用 CLIP 距离作为损失函数
\end{itemize}

例如重建一个推土机模型时，我们知道无论从哪个角度看，它都应该是``推土机''。因此，对于缺少参考图像的视角，我们可以要求渲染出的图像在 CLIP 嵌入空间中与``bulldozer''这个文本描述尽可能接近。

\subsubsection{零样本 3D 生成}

CLIP 引导的极限情况是：\textbf{完全没有任何参考图像}，仅凭一个文本描述就生成 3D 模型。这时候所有视角都使用 CLIP 损失：

$$
\theta^* = \arg\min_{\theta} \sum_{\pi \sim \mathcal{P}} d_{\text{CLIP}}(g(\theta, \pi),\ \text{``a bulldozer''})
$$

\begin{importantbox}{从重建到生成的范式转变}
当输入图像数量减少到零时，任务的性质发生了根本变化：不再是``3D 重建''，而是``3D 生成''。这是利用 2D 先验知识创造 3D 内容的第一次重要尝试。

代表性工作：DreamFields（2021）——仅用 CLIP 损失和 NeRF 实现了文本到 3D 的生成。
\end{importantbox}

\subsection{CLIP 引导的局限性}

虽然 CLIP 引导的 3D 生成开辟了新方向，但输出质量有限——生成的 3D 模型往往较为粗糙，细节不足。这是因为 CLIP 本质上是一个判别模型（discriminator），只能判断图像与文本是否匹配，但不能提供``图像应该长什么样''的精细指导。

\begin{warningbox}{CLIP 作为损失函数的局限}
CLIP 嵌入空间中的距离只能衡量语义级别的对齐，无法提供像素级别的生成指导。用 CLIP 损失优化时，模型可能会产生``语义正确但视觉不自然''的结果（adversarial examples 问题）。
\end{warningbox}

\subsection{核心问题：能否用扩散模型替代 CLIP？}

CLIP 是一个天然的判别器——输入图像和文本，输出匹配分数。但扩散模型是一个生成器——输入文本，输出图像。

\textbf{问题}：如何将一个生成模型当作判别器来使用？

如果我们使用的是 GAN，这很简单——GAN 本身就包含一个判别器。但扩散模型没有显式的判别器。Score Distillation Sampling（SDS）正是为解决这个问题而提出的。

\subsection{本章小结}

CLIP 引导的 3D 生成（如 DreamFields）证明了利用 2D 先验生成 3D 内容的可行性，但受限于 CLIP 只能提供语义级别的引导。下一步自然的想法是：用拥有更丰富像素级知识的扩散模型来替代 CLIP，从而获得更高质量的生成结果。

\newpage
